\documentclass[uplatex,dvipdfmx]{jlreq}
\usepackage[fleqn,tbtags]{mathtools}
\usepackage{jlreq-deluxe}
\usepackage[noalphabet]{pxchfon} % must be after otf package
\usepackage{stix2} %欧文＆数式フォント
\usepackage[fleqn,tbtags]{mathtools} % 数式関連 (w/ amsmath)
\usepackage{hira-stix} % ヒラギノフォント＆STIX2 フォント代替定義（Warning回避）

\usepackage{url}
\usepackage{float}

\usepackage{etoolbox}\pretocmd{\section}{\par\noindent\rule{\linewidth}{0.4pt}\par}{}{}
\usepackage[dvipdfmx,hidelinks]{hyperref}
\usepackage{pxjahyper}

\begin{document}

{\LARGE \textbf{作業ログ}}

報告者：佐藤 夏美（25G1062）

報告日：2026年5月26日

プロジェクト名：受動的環境介入によるメンタルヘルス支援：大学空間における自動演奏装置の提案

\section{進捗サマリー（要約）}
\begin{itemize}
  \item 全体進捗率：50\%（クラリネット音色設計）
  \item マイルストーン達成状況：進行中
  \item 主要成果：
  \begin{itemize}
    \item FFT解析完了・倍音比率の取得（実測クラリネット音源の奇数倍音優勢を確認）
    \item クラリネット音色作成（8倍音の加算合成・ローパス・ADSR）
  \end{itemize}
  \item 重大課題（影響度：高）：
  \begin{itemize}
    \item なし
  \end{itemize}
\end{itemize}

\section{作業ログ（詳細トラッキング）}
\begin{table}[H]
  \centering
  \begin{tabular}{p{2.2cm}p{1cm}p{1cm}p{1cm}p{1cm}p{2.2cm}p{2.2cm}}
\hline
作業項目 & 開始日 & 予定完了日 & 状態 & 進捗率 & 依存関係・ブロッカー & コメント \\ \hline \hline
音色調査（クラリネットの特徴） & 5/20 & 5/21 & 完了 & 100\% & - & 奇数倍音優勢を資料で確認 \\ \hline
FFT解析・倍音比率取得 & 5/21 & 5/22 & 完了 & 100\% & 実測音源 & 第8倍音まで実測 \\ \hline
clarinet1.pde作成 & 5/22 & 5/26 & 完了 & 100\% & FFTデータ & 8倍音・LPF・ADSR \\ \hline
  \end{tabular}
\end{table}

\section{課題管理}
\begin{table}[H]
  \centering
  \begin{tabular}{p{2.2cm}p{2.2cm}p{1.3cm}p{1.3cm}p{2.2cm}p{1cm}p{1.3cm}}
\hline
課題 & 発生原因 & 影響度 & 優先度 & 対策 & 期限 & 状態 \\ \hline \hline
機械音感 & ビブラートや楽器特有の音の鳴り方の不足 & 中 & 中 & 音域連動を調整，新たな要素の追加検討 & 6/2 & 対応中 \\ \hline
  \end{tabular}
\end{table}
※影響度＝プロジェクトへのインパクト　※優先度＝対応の緊急性

\section{AI利用状況と検証（再現性重視）}
\subsection{利用内容}
\begin{itemize}
  \item 利用目的：設計・実装
  \item 入力（プロンプト要旨）：
  \begin{itemize}
    \item 担当箇所の設計資料をもとに加算合成＋3系統ADSR＋音程連動ローパスでクラリネット音色設計を依頼
  \end{itemize}
  \item 出力概要：
  \begin{itemize}
    \item 倍音テーブルから \texttt{WavetableGenerator.gen10} で波形生成→\texttt{Oscil}→\texttt{MoogFilter}（LP）→\texttt{ADSR} と接続し，リードノイズ・ブレスノイズを別系統のADSRで重ねる設計案
  \end{itemize}
\end{itemize}

\subsection{検証方法}
※第三者が再現できるレベルで記述
\begin{itemize}
  \item 検証手法：
  \begin{itemize}
    \item FFT解析（実測との比較）
    \item 担当箇所の資料との照合
    \item 実機での実行と聴感確認
  \end{itemize}
  \item 参照資料：
  \begin{itemize}
    \item 実測に用いたクラリネット音源：MTG（Music Technology Group, Universitat Pompeu Fabra）, freesound.org（音源番号248359）\\
          \url{https://freesound.org/people/MTG/sounds/248359/}
    \item Processing 公式リファレンス：The Processing Foundation\\
          \url{https://processing.org/reference/}
    \item Minim 公式ドキュメント\\
          \url{https://code.compartmental.net/minim/}
  \end{itemize}
  \item 検証手順：
  \begin{itemize}
    \item 作成した音色が奇数倍音優勢かを確認する
    \item AI出力の関数構成を読み，想定通りになっているか確認する
    \item 再生し，クラリネットらしい音色になっているか聴いて確認する
  \end{itemize}
\end{itemize}

\subsection{ファクトチェック結果}
\begin{itemize}
  \item 一致点：
  \begin{itemize}
    \item FFT解析の結果，奇数倍音優勢の構造が確認できた
    \item 関数構成と処理経路（加算合成→LPF→ADSR）は担当箇所の設計資料および Minim・Processing 公式の記法と整合した
    \item 改善の余地はあるがクラリネットらしい音色を聴感より確認した
  \end{itemize}
  \item 不一致・疑義：
  \begin{itemize}
    \item AI提示のADSRやフィルタなどの値は実際の音と比べ調整の余地を感じた
  \end{itemize}
  \item 最終判断：
  \begin{itemize}
    \item 一部採用．関数構成・処理経路はAI案を採用し，ADSRやフィルタなどの値は調節した
  \end{itemize}
  \item 根拠：
  \begin{itemize}
    \item FFT解析および聴感による結果
  \end{itemize}
\end{itemize}

\section{次回計画}
\begin{itemize}
  \item 目標：
  \begin{itemize}
    \item 音色をさらに本物に近づける
  \end{itemize}
  \item 達成基準：
  \begin{itemize}
    \item FFT解析結果をより実音源に近づけること
    \item 実音源との比較でより自然なクラリネット音色になっていること
  \end{itemize}
  \item 期限：
  \begin{itemize}
    \item 2026年6月2日
  \end{itemize}
\end{itemize}

\section{付録}
\begin{itemize}
  \item 添付資料：
  \begin{itemize}
    \item ソースコード：\url{https://github.com/7nanasi/hackathon}
  \end{itemize}
  \item 添付範囲：差分（今回着手分）
\end{itemize}

\end{document}